El mes de novembre del 2021, la NASA va llançar la missió DART (Double Asteroid Redirection Test). Aquesta missió té per objectiu canviar l'òrbita de Dimorphos, un petit asteroide que orbita al voltant de Didymos, que és un asteroide més gran.

\figura[0.6]{fig1.png}
\begin{center}
\begin{minipage}{0.6\linewidth}\small
Esquema de la missió DART.\\
\textsc{Nota}: La imatge no està escalada, les proporcions entre les òrbites i els objectes no es corresponen a les escales reals.\\
\textsc{Font}: NASA.
\end{minipage}
\end{center}

\begin{apartat}{1,25}
A partir de la llei de la gravitació universal, trobeu l'expressió de la intensitat del camp gravitatori que crea un objecte astronòmic esfèric de massa $M$ i radi $R$ a la seva superfície. El diàmetre de Didymos és de 781~m i la seva densitat és de $2\,146\ \mathrm{kg/m^3}$. Calculeu el valor de la intensitat del camp gravitatori que crea Didymos a la seva superfície. Si Dimorphos té una massa de $4,42\times 10^{10}\ \mathrm{kg}$ i el radi orbital mitjà (distància entre els centres dels dos objectes) és d'1,12~km, calculeu el mòdul de la força gravitatòria mitjana entre Didymos i Dimorphos.
\end{apartat}

\begin{apartat}{1,25}
L'objectiu de la missió DART és colpejar Dimorphos, de tal manera que orbiti en una nova òrbita de radi menor, com s'indica en la figura anterior. Deduïu, a partir de principis fonamentals, l'expressió de la velocitat orbital d'un satèl·lit en funció del radi de l'òrbita. Argumenteu si Dimorphos orbitarà a més velocitat a la nova òrbita o a l'òrbita original.
\end{apartat}

\dades{%
  $G = 6,67\times 10^{-11}\ \mathrm{N\,m^2\,kg^{-2}}$.}
\nota{Considereu que les òrbites són circulars i que els dos asteroides són esfèrics.}
